\documentclass[10pt,a4paper]{amsart}
\usepackage{amsmath,amssymb,amsthm,graphicx}
\usepackage{verbatim}
\newenvironment{reference}{\comment}{\endcomment}
%\newenvironment{reference}{}{}
\newenvironment{slogan}{\comment}{\endcomment}
\newenvironment{history}{\comment}{\endcomment}

% For commutative diagrams we use Xy-pic
\usepackage[all]{xy}

% We use 2cell for 2-commutative diagrams.
\xyoption{2cell}
% Theorem environments.
%
\theoremstyle{plain}
\newtheorem{theorem}[subsection]{Theorem}
\newtheorem{proposition}[subsection]{Proposition}
\newtheorem{lemma}[subsection]{Lemma}

\theoremstyle{definition}
\newtheorem{definition}[subsection]{Definition}
\newtheorem{example}[subsection]{Example}
\newtheorem{exercise}[subsection]{Exercise}
\newtheorem{situation}[subsection]{Situation}

\theoremstyle{remark}
\newtheorem{remark}[subsection]{Remark}
\newtheorem{remarks}[subsection]{Remarks}

\numberwithin{equation}{subsection}

% Macros
%
\def\lim{\mathop{\mathrm{lim}}\nolimits}
\def\colim{\mathop{\mathrm{colim}}\nolimits}
\def\Spec{\mathop{\mathrm{Spec}}}
\def\Hom{\mathop{\mathrm{Hom}}\nolimits}
\def\Ext{\mathop{\mathrm{Ext}}\nolimits}
\def\SheafHom{\mathop{\mathcal{H}\!\mathit{om}}\nolimits}
\def\SheafExt{\mathop{\mathcal{E}\!\mathit{xt}}\nolimits}
\def\Sch{\mathit{Sch}}
\def\Mor{\mathop{\mathrm{Mor}}\nolimits}
\def\Ob{\mathop{\mathrm{Ob}}\nolimits}
\def\Sh{\mathop{\mathit{Sh}}\nolimits}
\def\NL{\mathop{N\!L}\nolimits}
\def\CH{\mathop{\mathrm{CH}}\nolimits}
\def\proetale{{pro\text{-}\acute{e}tale}}
\def\etale{{\acute{e}tale}}
\def\QCoh{\mathit{QCoh}}
\def\Ker{\mathop{\mathrm{Ker}}}
\def\Im{\mathop{\mathrm{Im}}}
\def\Coker{\mathop{\mathrm{Coker}}}
\def\Coim{\mathop{\mathrm{Coim}}}

% Boxtimes
%
\DeclareMathSymbol{\boxtimes}{\mathbin}{AMSa}{"02}

%
% Macros for moduli stacks/spaces
%
\def\QCohstack{\mathcal{QC}\!\mathit{oh}}
\def\Cohstack{\mathcal{C}\!\mathit{oh}}
\def\Spacesstack{\mathcal{S}\!\mathit{paces}}
\def\Quotfunctor{\mathrm{Quot}}
\def\Hilbfunctor{\mathrm{Hilb}}
\def\Curvesstack{\mathcal{C}\!\mathit{urves}}
\def\Polarizedstack{\mathcal{P}\!\mathit{olarized}}
\def\Complexesstack{\mathcal{C}\!\mathit{omplexes}}
% \Pic is the operator that assigns to X its picard group, usage \Pic(X)
% \Picardstack_{X/B} denotes the Picard stack of X over B
% \Picardfunctor_{X/B} denotes the Picard functor of X over B
\def\Pic{\mathop{\mathrm{Pic}}\nolimits}
\def\Picardstack{\mathcal{P}\!\mathit{ic}}
\def\Picardfunctor{\mathrm{Pic}}
\def\Deformationcategory{\mathcal{D}\!\mathit{ef}}
\usepackage{cleveref}
\crefname{originalsection}{original section}{original sections}
\Crefname{originalsection}{Original section}{Original sections}
\makeatletter
\newlabel{constructions-section-projective-space}{{01ND}{source}{}{sourcepointer-1}{}}
\newlabel{constructions-section-projective-space@cref}{{[originalsection][0][]01ND}{[1][0][]source}{}{}{}}
\makeatother
\crefname{originallemma}{original lemma}{original lemmas}
\Crefname{originallemma}{Original lemma}{Original lemmas}
\makeatletter
\newlabel{descent-lemma-descending-property-closed-immersion}{{02L6}{source}{}{sourcepointer-2}{}}
\newlabel{descent-lemma-descending-property-closed-immersion@cref}{{[originallemma][0][]02L6}{[1][0][]source}{}{}{}}
\makeatother
\crefname{originalreference}{original reference}{original references}
\Crefname{originalreference}{Original reference}{Original references}
\makeatletter
\newlabel{lemma-automatic}{{0BY5}{source}{}{sourcepointer-3}{}}
\newlabel{lemma-automatic@cref}{{[originalreference][0][]0BY5}{[1][0][]source}{}{}{}}
\makeatother
\crefname{originalreference}{original reference}{original references}
\Crefname{originalreference}{Original reference}{Original references}
\makeatletter
\newlabel{lemma-criterion-very-ample}{{0E8V}{source}{}{sourcepointer-4}{}}
\newlabel{lemma-criterion-very-ample@cref}{{[originalreference][0][]0E8V}{[1][0][]source}{}{}{}}
\makeatother
\title{Very ampleness of the square of a globally generated invertible module with vanishing H\textasciicircum{}1}
\author{The Stacks Project Authors}
\date{stacks-project:curves.tex:local-snapshot-20261012; GNU FDL 1.2 or later}
\begin{document}
\maketitle
\noindent\textit{Editorial note.} All delivered mathematical text is the author's own, copied from the pinned Stacks Project chapter file sources/stacks-curves/pinned-original-source.tex; no statement or proof step was rewritten, reordered or removed.\par\medskip
\noindent\textit{Editorial note.} The counted proof is the author's complete argument: the proof twice delegates to the proof of the registered sibling \textbackslash{}label\{lemma-criterion-very-ample\} (source lines 1636-1637: 'all of the statements given in the proof of Lemma criterion-very-ample hold with this s'; lines 1650-1652: 'Arguing as in the proof of Lemma criterion-very-ample we find that S\textasciicircum{}\{(2)\} is generated in degree 1'), and ends with 'Some details omitted.' (line 1654); the essential skeleton - the descent reduction, the regular section s via lemma-automatic, the extra generator t giving surjectivity of mu\_n, and the generation of S in degrees 0, 1, 2 - is present in the proof; the delegated statements are available in the registered sibling card.\par\medskip
\noindent\textit{Editorial note.} The label lemma-criterion-very-ample-bis names the counted statement, and the proof environment immediately following it in the source is delivered verbatim as its proof.\par\medskip
\noindent\textit{Editorial note.} The counted statement is carried by the author's own lemma environment in the chapter source.\par\medskip
\noindent\textit{Allowed standard prerequisites (editorial).}\begin{itemize}
\item The morphism defined by a basis of global sections. Constructions, Section \textbackslash{}label\{constructions-section-projective-space\}. Proof lines 1615-1617: global generation of L\textasciicircum{}\{\textbackslash{}otimes 2\} gives the morphism phi to P\textasciicircum{}n\_k associated to a chosen basis s\_0, ..., s\_n of H\textasciicircum{}0(X, L\textasciicircum{}\{\textbackslash{}otimes 2\}).
\item Descending the closed immersion property. Descent, Lemma \textbackslash{}label\{descent-lemma-descending-property-closed-immersion\}. Proof lines 1617-1621: to check that the morphism is a closed immersion one may replace k by its algebraic closure, so assume k algebraically closed.
\item Automatic Cohen-Macaulay property of the curve. Same chapter, \textbackslash{}label\{lemma-automatic\}. Proof lines 1629-1633: X is Cohen-Macaulay, hence the invertible module L has no embedded associated points and the section s nonvanishing at all generic points is a regular section.
\item The sibling very ampleness criterion. Same chapter, \textbackslash{}label\{lemma-criterion-very-ample\}. Proof lines 1636-1637 and 1650-1652: all statements in the proof of this registered sibling hold with the chosen s, and the final generation-in-degree-1 argument for S\textasciicircum{}\{(2)\} is the same as in its proof.
\end{itemize}\medskip
\ifdefined\hypertarget\hypertarget{sourcepointer-1}{}\fi
\noindent\textit{Source reference.} Original section 01ND: sources/stacks-chow/tags.txt, source lines 2154--2154. Prerequisite pointer invoked at proof line(s) 1614; it names a standard fact of a sibling chapter that the counted argument uses, and every step of the counted proof built on it is retained verbatim in the delivered body.\par\medskip
\ifdefined\hypertarget\hypertarget{sourcepointer-2}{}\fi
\noindent\textit{Source reference.} Original lemma 02L6: sources/stacks-chow/tags.txt, source lines 3397--3397. Prerequisite pointer invoked at proof line(s) 1617; it names a standard fact of a sibling chapter that the counted argument uses, and every step of the counted proof built on it is retained verbatim in the delivered body.\par\medskip
\ifdefined\hypertarget\hypertarget{sourcepointer-3}{}\fi
\noindent\textit{Source reference.} Original reference 0BY5: sources/stacks-curves/pinned-original-source.tex, source lines 1277--1277. Prerequisite pointer invoked at proof line(s) 1628; it names a standard fact of the same chapter that the counted argument uses, and every step of the counted proof built on it is retained verbatim in the delivered body.\par\medskip
\ifdefined\hypertarget\hypertarget{sourcepointer-4}{}\fi
\noindent\textit{Source reference.} Original reference 0E8V: sources/stacks-curves/pinned-original-source.tex, source lines 1506--1506. Prerequisite pointer invoked at proof line(s) 1633; it names a standard fact of the same chapter that the counted argument uses, and every step of the counted proof built on it is retained verbatim in the delivered body.\par\medskip
\begin{lemma}
\label{lemma-criterion-very-ample-bis}
Let $k$ be a field. Let $X$ be a proper scheme over $k$ having dimension $1$
and $H^0(X, \mathcal{O}_X) = k$. Let $\mathcal{L}$ be an invertible
$\mathcal{O}_X$-module. Assume
\begin{enumerate}
\item $\mathcal{L}$ is globally generated,
\item $H^1(X, \mathcal{L}) = 0$, and
\item $\mathcal{L}$ is ample.
\end{enumerate}
Then $\mathcal{L}^{\otimes 2}$ is very ample on $X$ over $k$.
\end{lemma}

\begin{proof}
Choose basis $s_0, \ldots, s_n$ of $H^0(X, \mathcal{L}^{\otimes 2})$
over $k$. By property (1) we see that $\mathcal{L}^{\otimes 2}$
is globally generated and we get a morphism
$$
\varphi_{\mathcal{L}^{\otimes 2}, (s_0, \ldots, s_n)} :
X \longrightarrow \mathbf{P}^n_k
$$
See Constructions, Section \ref{constructions-section-projective-space}.
The lemma asserts that this morphism is a closed immersion.
To check this we may replace $k$ by its algebraic closure, see
Descent, Lemma \ref{descent-lemma-descending-property-closed-immersion}.
Thus we may assume $k$ is algebraically closed.

\medskip\noindent
Assume $k$ is algebraically closed. For each generic point $\eta_i \in X$
let $V_i \subset H^0(X, \mathcal{L})$ be the $k$-subvector space of
sections vanishing at $\eta_i$. Since $\mathcal{L}$ is globally generated,
we see that $V_i \not = H^0(X, \mathcal{L})$. Since $X$ has only a
finite number of irreducible components and $k$ is infinite, we can find
$s \in H^0(X, \mathcal{L})$ nonvanishing at $\eta_i$ for all $i$.
Then $s$ is a regular section of $\mathcal{L}$ (because $X$ is
Cohen-Macaulay by Lemma \ref{lemma-automatic} and hence $\mathcal{L}$
has no embedded associated points).

\medskip\noindent
In particular, all of the statements given in the proof of
Lemma \ref{lemma-criterion-very-ample} hold with this $s$.
Moreover, as $\mathcal{L}$ is globally generated, we can find
a global section $t \in H^0(X, \mathcal{L})$ such that
$t|_Z$ is nonvanishing (argue as above using the finite number
of points of $Z$). Then in the proof of Lemma \ref{lemma-criterion-very-ample}
we can use $t$ to see that additionally the multiplication map
$$
\mu_n :
H^0(X, \mathcal{L}) \otimes_k H^0(X, \mathcal{L}^{\otimes 2})
\longrightarrow
H^0(X, \mathcal{L}^{\otimes 3})
$$
is surjective. Thus
$$
S = \bigoplus\nolimits_{n \geq 0} H^0(X, \mathcal{L}^{\otimes n})
$$
is generated in degrees $0, 1, 2$ over $k$. Arguing as in the
proof of Lemma \ref{lemma-criterion-very-ample} we find that
$S^{(2)} = \bigoplus_{n} S_{2n}$ is generated in degree $1$.
This means that $\mathcal{L}^{\otimes 2}$ is very ample as desired.
Some details omitted.
\end{proof}

\end{document}
